\chapter{Architettura del control plane}
\label{cap:04-architettura-control-plane}

\section{Vista d'insieme}

Il control plane di \nf{} \`e un'applicazione Spring Boot che espone due porte:
la \code{8080} per l'API di gestione e invocazione, e la \code{8081} per gli
endpoint di management --- salute e metriche Prometheus. La separazione non \`e
cosmetica: consente di esporre l'API al traffico applicativo senza esporre gli
endpoint diagnostici, e soprattutto fa s\`i che lo scraping delle metriche non
condivida n\'e le code n\'e i limiti di ammissione del traffico che sta
misurando. Uno strumento di misura che si mette in coda dietro il carico che
osserva riporta la propria attesa insieme a quella del sistema.

La \cref{fig:architettura} d\`a la vista d'insieme. Il riquadro esterno \`e il
processo del control plane; al suo interno la fascia superiore contiene i quattro
componenti del nucleo, la fascia inferiore tratteggiata i moduli opzionali. Fuori
dal riquadro, in basso, i tre tipi di destinazione verso cui un'invocazione pu\`o
essere dispacciata.

\begin{figure}[htbp]
  \centering
  \resizebox{0.95\textwidth}{!}{% nanofaas system architecture (icon / "card" style).
% Grouping containers live on the background layer and are declared AFTER every
% blur-shadow node (shadows.blur manipulates layers; a shadowed node emitted
% after background content suppresses it). Keep that ordering when editing.
\documentclass[border=8pt]{standalone}
\usepackage{nanofaas-figs}
\begin{document}
\begin{tikzpicture}[nf]

  % ================= core components (shadowed cards) =================
  \node[iconcore, minimum width=44mm, minimum height=27mm] (registry)
    {{\color{nfnavy}\large\ic{book}}\ \ \textbf{Function Registry}\\[5pt]
     {\footnotesize\color{nfslate}in-memory}};
  \node[iconcore, right=6mm of registry, minimum width=48mm, minimum height=27mm] (invoc)
    {{\color{nfnavy}\large\ic{server}}\ \ \textbf{Invocation Service}\\[5pt]
     {\footnotesize\color{nfslate}sync/async $\cdot$ retries}\\
     {\footnotesize\color{nfslate}idempotency $\cdot$ rate limit}};
  \node[iconcore, right=6mm of invoc, minimum width=46mm, minimum height=27mm] (dispatch)
    {{\color{nfnavy}\large\ic{sitemap}}\ \ \textbf{Pool Dispatcher}\\[5pt]
     {\footnotesize\color{nfslate}LOCAL / POOL /}\\
     {\footnotesize\color{nfslate}DEPLOYMENT}};
  \node[iconcore, right=6mm of dispatch, minimum width=44mm, minimum height=27mm] (store)
    {{\color{nfnavy}\large\ic{database}}\ \ \textbf{Execution Store}\\[5pt]
     {\footnotesize\color{nfslate}lifecycle $\cdot$ TTL eviction}};

  % ================= optional modules (shadowed cards in a grid) =================
  \matrix (mods) [matrix of nodes, anchor=north,
      nodes={iconmod, minimum width=59mm, anchor=center},
      column sep=5mm, row sep=3mm]
      at ($(registry.south)!0.5!(store.south)+(0,-17mm)$) {
    |[iconmod]| {\color{nfgreen}\ic{envelope}}~~async-queue &
    |[iconmod]| {\color{nfgreen}\ic{sync-alt}}~~sync-queue &
    |[iconmod]| {\color{nfgreen}\ic{chart-bar}}~~autoscaler \\
    |[iconmod]| {\color{nfgreen}\ic{cloud-upload-alt}}~~offload &
    |[iconmod]| {\color{nfgreen}\ic{cog}}~~runtime-config &
    |[iconmod]| {\color{nfgreen}\ic{shield-alt}}~~image-validator \\
    |[iconmod]| {\color{nfgreen}\ic{dharmachakra}}~~k8s-deployment-provider &
    |[iconmod]| {\color{nfgreen}\ic{docker}}~~container-deployment-provider &
    |[iconmod]| {\color{nfgreen}\ic{tag}}~~build-metadata \\
  };

  % ================= execution targets (shadowed cards) =================
  \node[iconback, minimum width=50mm, minimum height=21mm] (clocal)
    at ($(mods.south)+(0,-22mm)$)
    {{\color{nfamber}\large\ic{docker}}\ \ \textbf{Local containers}\\[4pt]
     {\footnotesize\color{nfslate}container-local $\cdot$ Docker CLI}};
  \node[iconback, left=9mm of clocal, minimum width=48mm, minimum height=21mm] (k8s)
    {{\color{nfamber}\large\ic{dharmachakra}}\ \ \textbf{Kubernetes}\\[4pt]
     {\footnotesize\color{nfslate}Deployment + Service}};
  \node[iconback, right=9mm of clocal, minimum width=42mm, minimum height=21mm] (pool)
    {{\color{nfamber}\large\ic{fire}}\ \ \textbf{Warm pod pool}\\[4pt]
     {\footnotesize\color{nfslate}OpenWhisk-style}};

  % ================= client =================
  \node[iconcore, minimum width=104mm, minimum height=23mm] (client)
    at ($(registry.north)!0.5!(store.north)+(0,44mm)$)
    {{\color{nfnavy}\LARGE\ic{users}}\ \ \textbf{\large Client / SDK}\quad
     \begin{tabular}[c]{@{}l@{}}
       {\footnotesize\ttfamily POST /v1/functions/\{name\}:invoke~(sync)}\\
       {\footnotesize\ttfamily POST /v1/functions/\{name\}:enqueue~(async)}
     \end{tabular}};

  % strip reserving room for the banner at the top of the control plane
  \node[inner sep=0, minimum height=10mm, anchor=south] (topstrip)
    at ($(registry.north)!0.5!(store.north)+(0,3mm)$) {};

  % ================= grouping containers (background, AFTER all shadowed nodes) ===
  \begin{scope}[on background layer]
    \node[container, draw=nfblue, line width=1.6pt, fill=nfblue!4,
      fit=(topstrip)(registry)(store)(mods)] (cp) {};
    \node[container, draw=nfgreen, line width=1.1pt, dash pattern=on 3pt off 2.4pt,
      fill=nfgreen!7, fit=(mods)] (modbox) {};
  \end{scope}

  % ================= banner + labels (no shadow) =================
  \node[banner, anchor=west] at ($(registry.north west)+(-1mm,8mm)$)
    {NanoFaaS Control Plane\, {\normalsize\mdseries(Java / Spring Boot)}};
  \node[font=\sffamily\small, text=nfnavy, anchor=east]
    at ($(store.north east)+(1mm,8mm)$) {\texttt{:8080} API \; $\cdot$ \; \texttt{:8081} metrics};
  \node[clabel, text=nfgreen, anchor=south west] at ($(mods.north west)+(-1mm,4mm)$)
    {\textbf{Optional modules}\, {\mdseries\footnotesize(loaded via \texttt{ControlPlaneModule} SPI)}};

  % ================= edges =================
  \draw[flow] (client.south -| invoc) -- (invoc |- cp.north);

  \coordinate (hub) at ($(cp.south)+(0,-9mm)$);
  \draw[flowc] (cp.south) -- (hub)
    node[tag, midway, right=1mm, fill=white, inner sep=1.5pt]
      {Dispatch to execution targets $\cdot$ LOCAL / POOL / DEPLOYMENT};
  \draw[flowc] (hub) -| (k8s.north);
  \draw[flowc] (hub) -- (clocal.north);
  \draw[flowc] (hub) -| (pool.north);

\end{tikzpicture}
\end{document}
}
  \caption{Architettura di \nf{}: control plane a processo singolo, moduli
  opzionali caricati via SPI, backend di esecuzione sostituibili.}
  \label{fig:architettura}
\end{figure}

I quattro componenti del nucleo si dividono le responsabilit\`a nel modo
seguente. Il primo e il quarto vivono in moduli diversi da quello del control
plane (\cref{cap:08-sistema-moduli}): lo store, la macchina a stati dei
tentativi e il motore di scheduling stanno in \code{execution-runtime}, mentre
i contratti con cui i moduli si innestano nel nucleo (le interfacce
\code{ControlPlaneModule}, dei provider e dei punti di estensione) stanno nella
libreria \code{control-plane-spi}.

\begin{description}[style=nextline, leftmargin=0pt]
  \item[\code{FunctionRegistry}] custodisce le specifiche delle funzioni
  registrate. \`E uno snapshot immutabile sostituito per intero a ogni scrittura
  e reso durevole su file; il \cref{cap:06-nucleo} ne descrive la
  persistenza, la risoluzione dei default e l'ordine di registrazione e
  rimozione.

  \item[\code{InvocationService}] orchestra le invocazioni: risolve la
  specifica, crea o riusa un record di esecuzione, e affida il resto al
  coordinatore reattivo. Il limite di frequenza non \`e pi\`u qui: sta in un
  filtro che precede la lettura del corpo (\cref{sec:04-percorso}).

  \item[\code{DispatcherRouter}] sceglie il dispatcher concreto in base alla
  modalit\`a di esecuzione della funzione (\cref{cap:07-modalita-esecuzione}).

  \item[\code{ExecutionStore}] mantiene le esecuzioni vive e gli esiti di quelle
  concluse, in due cache con scadenza e limite propri; la transizione terminale
  \`e di \code{ExecutionLifecycle}.
\end{description}

\section{Perch\'e reattivo}

Il control plane usa Spring WebFlux~\cite{springwebflux} su Netty, non lo stack
servlet. La scelta \`e determinata dal profilo di carico che il componente
sostiene.

Il control plane non calcola: attende. Per ogni invocazione sincrona, il lavoro
che gli compete --- validazione, risoluzione della specifica, creazione del
record --- occupa qualche decina di microsecondi; poi il processo resta in
attesa della risposta della funzione, che dura da qualche millisecondo a
decine di secondi. Con un modello a thread-per-richiesta, sostenere $N$
invocazioni contemporanee richiede $N$ thread bloccati, ciascuno con il proprio
stack. A $N$ nell'ordine delle migliaia il costo di memoria degli stack diventa
comparabile all'intera impronta del processo, e il costo di context switch
inizia a comparire nella latenza misurata --- ossia nella grandezza che gli
esperimenti di questo documento devono osservare pulita.

Con il modello reattivo la richiesta in attesa \`e un oggetto in heap, non un
thread. Il costo per richiesta in volo scende di due ordini di grandezza, e il
numero di thread resta legato al numero di core anzich\'e al numero di
richieste.

\subsection{Il prezzo: la disciplina del non bloccare}

Il modello reattivo impone un vincolo che la programmazione bloccante non ha:
nessuna operazione che possa attendere deve essere eseguita su un thread
dell'event loop. Un event loop bloccato non rallenta una richiesta, le blocca
tutte.

Nel codice di \nf{} questo vincolo compare in due punti espliciti, entrambi
commentati come tali. Il primo \`e la preparazione dell'invocazione sincrona, e
la sua forma \`e cambiata in modo istruttivo:

\begin{lstlisting}[language=Java, caption={Confinamento condizionale della fase
preparatoria (\code{InvocationService}).}]
Mono<Prepared> prepared = Mono.fromCallable(() -> {
    FunctionSpec spec = functionService.get(functionName)
            .orElseThrow(FunctionNotFoundException::new);
    refuseEarlyIfQueueFull(functionName, spec, idempotencyKey, InvocationKind.SYNC);
    return new Prepared(spec, executionFactory.createOrReuseExecution(
            functionName, spec, request, idempotencyKey, traceId, InvocationKind.SYNC));
});
// Hop only when there is something to hop for. ...
if (idempotencyKey != null && !idempotencyKey.isBlank()) {
    prepared = prepared.subscribeOn(Schedulers.boundedElastic());
}
return prepared.flatMap(p -> reactiveCoordinator.invoke(
        p.lookup(), p.spec(), timeoutOverrideMs, offloadContext));
\end{lstlisting}

La risoluzione dell'idempotenza usa un ciclo di \emph{compare-and-set} su una
cache concorrente, che sotto contesa pu\`o attendere. Una versione precedente
confinava sullo scheduler elastico \emph{ogni} richiesta; il commento attuale
spiega perch\'e non pi\`u: \code{createOrReuseExecution} sosta solo dentro quel
ciclo e, senza chiave, ritorna al primo ramo senza toccare alcun lock, e il resto
del blocco \`e un controllo di coda e una ricerca in mappa. Il passaggio a un
altro thread costa 6,1~$\mu$s misurati contro 0,085~$\mu$s di lavoro, settanta
volte ci\`o che protegge, e lo pagavano anche le richieste che sarebbero state
rifiutate, per le quali decidere un \code{429} non dovrebbe richiedere un secondo
thread. Oggi il salto avviene \emph{solo} se la richiesta porta una chiave.

Il secondo punto \`e l'intero percorso asincrono (\code{:enqueue}), che il
controller esegue sempre sullo scheduler elastico: l'accodamento pu\`o contendere
sullo stato della coda della funzione, e non deve toccare l'event loop.

\subsection{Il completamento come futuro condiviso}

Il punto architetturalmente pi\`u interessante del percorso sincrono \`e come
il control plane attende l'esito. Ogni \code{ExecutionRecord} possiede un
\code{CompletableFuture} che verr\`a completato da chiunque porti l'esito ---
il dispatcher che riceve la risposta HTTP, oppure la callback che la funzione
invia all'endpoint interno \code{/v1/internal/executions/\{id\}:complete}. Il
percorso sincrono si limita a sottoscrivere quel futuro:

\begin{lstlisting}[language=Java, caption={Attesa dell'esito nel
\code{ReactiveInvocationCoordinator}.}]
// suppressCancel=true: a single subscriber's timeout/disconnect must not cancel
// the shared completion future other idempotent waiters depend on.
return Mono.fromFuture(executionRecord.completion(), true)
        .onErrorMap(CompletionException.class,
                ex -> ex.getCause() != null ? ex.getCause() : ex)
        .timeout(Duration.ofMillis(timeoutMs))
        ...
\end{lstlisting}

Il parametro \code{suppressCancel} merita attenzione perch\'e codifica un
requisito che nasce dall'idempotenza. Se due chiamanti inviano la stessa chiave
di idempotenza, entrambi finiscono ad attendere lo \emph{stesso} futuro. Se il
primo scade o chiude la connessione, la cancellazione del suo \code{Mono} non
deve propagarsi al futuro condiviso, o il secondo chiamante osserverebbe un
fallimento causato dall'abbandono di qualcun altro. La soppressione della
cancellazione \`e ci\`o che rende l'idempotenza corretta anche in presenza di
disconnessioni.

\section{Il percorso di una richiesta}
\label{sec:04-percorso}

La \cref{fig:flusso} contrappone i due percorsi. Vale la pena percorrerli in
prosa, perch\'e l'ordine delle decisioni \`e esso stesso un elemento di
progetto.

\begin{figure}[htbp]
  \centering
  \resizebox{0.95\textwidth}{!}{% nanofaas invocation flow: synchronous (:invoke) and asynchronous (:enqueue)
% paths as a UML-style sequence diagram with an alt combined fragment.
\documentclass[border=8pt]{standalone}
\usepackage{nanofaas-figs}
\begin{document}
\begin{tikzpicture}[nf]

  % ---- lifeline x-coordinates (cm) ----
  \def\cx{0}\def\px{4.0}\def\qx{8.2}\def\dx{11.5}\def\rx{15.2}
  \def\ybot{-10.6}

  % ---- participant heads ----
  \node[client, minimum width=22mm] (cH) at (\cx,0) {Client / SDK};
  \node[core,   minimum width=34mm] (pH) at (\px,0)
    {Control Plane\\{\scriptsize validate $\cdot$ rate limit $\cdot$ ExecutionStore}};
  \node[module, minimum width=30mm] (qH) at (\qx,0)
    {Admission queue\\{\scriptsize sync-queue / async-queue}};
  \node[core,   minimum width=24mm] (dH) at (\dx,0) {DispatcherRouter};
  \node[backend,minimum width=26mm] (rH) at (\rx,0)
    {Function runtime\\{\scriptsize LOCAL $\cdot$ EXTERNAL $\cdot$ DEPLOYMENT}};

  % ---- lifelines ----
  \foreach \n in {cH,pH,qH,dH,rH}{\draw[lifeline] (\n.south) -- ($(\n.south)-(0,10.1)$);}

  % ---- message helpers ----
  \newcommand{\msg}[4]{\draw[flow] (#1,#3) -- (#2,#3)
      node[tag, midway, above=0.2mm]{#4};}
  \newcommand{\rmsg}[4]{\draw[ret] (#1,#3) -- (#2,#3)
      node[tagc, midway, above=0.2mm]{#4};}

  % ================= alt combined fragment frame =================
  \coordinate (fnw) at (-1.35,-0.7);
  \coordinate (fse) at (16.15,\ybot+0.35);
  \draw[draw=nfslate, line width=0.7pt] (fnw) rectangle (fse);
  \node[anchor=north west, fill=nfslatebg, draw=nfslate, line width=0.5pt,
        font=\sffamily\scriptsize\bfseries, inner sep=2pt, text=nfink] at (fnw) {alt};
  \node[anchor=west, tagc, font=\sffamily\scriptsize\itshape]
    at ($(fnw)+(0.95,-0.16)$) {[\,synchronous $\cdot$ \texttt{:invoke}\,]};

  % ---------- synchronous operand ----------
  \msg{\cx}{\px}{-1.35}{\texttt{POST \dots:invoke}}
  % self-activity on the control plane
  \draw[flow] ($(\px,-1.95)$) -- ++(0.9,0) -- ++(0,-0.5) -- ++(-0.9,0);
  \node[tag, anchor=west] at ($(\px+0.95,-2.2)$) {validate $\cdot$ rate limit $\cdot$ create exec state};
  \msg{\px}{\qx}{-2.9}{admit \scriptsize(else \texttt{429})}
  \msg{\qx}{\dx}{-3.5}{dispatch}
  \msg{\dx}{\rx}{-4.1}{forward \texttt{/invoke}}
  \rmsg{\rx}{\dx}{-4.6}{result}
  \rmsg{\dx}{\px}{-5.0}{update store}
  \rmsg{\px}{\cx}{-5.4}{\texttt{200} + body}

  % ---------- operand divider ----------
  \draw[dashed, draw=nfslate] (-1.35,-5.9) -- (16.15,-5.9);
  \node[anchor=west, tagc, font=\sffamily\scriptsize\itshape]
    at (-1.30,-6.06) {[\,asynchronous $\cdot$ \texttt{:enqueue}\,]};

  % ---------- asynchronous operand ----------
  \msg{\cx}{\px}{-6.6}{\texttt{POST \dots:enqueue}}
  \msg{\px}{\qx}{-7.15}{enqueue \scriptsize(needs async-queue; else \texttt{501})}
  \rmsg{\px}{\cx}{-7.6}{\texttt{202} + execution id}

  \draw[dotted, draw=nfslate] (-1.0,-8.0) -- (15.2,-8.0);
  \node[fill=white, tagc, font=\sffamily\scriptsize\itshape] at (7.1,-8.0)
    {scheduler drains the queue asynchronously};

  \msg{\qx}{\dx}{-8.55}{dispatch}
  \msg{\dx}{\rx}{-9.1}{forward \texttt{/invoke}}
  \rmsg{\rx}{\dx}{-9.55}{result}
  \rmsg{\dx}{\px}{-9.95}{update store}
  \msg{\px}{\cx}{-10.35}{callback \scriptsize POST \texttt{callbackUrl}}

\end{tikzpicture}
\end{document}
}
  \caption{Invocazione sincrona (\code{:invoke}) e asincrona (\code{:enqueue})
  a confronto.}
  \label{fig:flusso}
\end{figure}

\subsection{Percorso sincrono}

\begin{enumerate}[leftmargin=*]
  \item Due filtri precedono il controller. Il primo, con precedenza massima,
  applica il confine del corpo (\code{413}, \cref{sec:06-capacita}); il
  secondo \`e il limitatore di frequenza globale, ristretto ai suffissi
  \code{:invoke} e \code{:enqueue}: una richiesta respinta risponde \code{429} con
  \code{Retry-After: 1} senza aver decodificato il JSON n\'e cercato la funzione,
  e il suo corpo viene comunque svuotato per non impedire il riuso della
  connessione.

  \item Il controller estrae gli header dedicati (chiave di idempotenza, trace
  id, override del timeout, marcatore di salto di offload) e
  \textbf{ricostruisce} la richiesta con gli header effettivi del chiamante.

  \item Il nome viene risolto nel registro; se assente, \code{404}. Se la coda
  della funzione \`e certamente piena e nulla pu\`o deviare la richiesta altrove, il
  rifiuto avviene qui, prima di costruire un'esecuzione destinata a essere
  abbandonata --- salvo quando \`e presente una chiave, perch\'e la replica di
  un'esecuzione esistente non deve essere respinta (nel commento, 590 rifiuti al
  secondo contro 299 dispatch).

  \item Si crea o si riusa il record di esecuzione, dopo aver riservato le
  quote di ammissione (\cref{sec:06-capacita}). Se la chiave di idempotenza
  corrisponde a un'esecuzione gi\`a conclusa, l'esito viene riprodotto
  immediatamente senza toccare la funzione (\emph{replay}); se l'esito \`e stato
  espulso per capacit\`a, la risposta \`e \code{410}.

  \item \textbf{Ammissione.} \`E il punto in cui il comportamento dipende dai
  moduli presenti, e la logica \`e una catena di ripiego ordinata: se il modulo
  di offload \`e attivo e la politica della funzione impone l'offload eager, la
  richiesta parte verso l'istanza remota; altrimenti si tenta l'ammissione
  locale, che a sua volta usa la coda sincrona se presente, la coda per funzione se
  presente, e in ultima istanza dispaccia direttamente, acquisendo un lease di
  capacit\`a: senza spazio la risposta \`e \code{429}, mai una coda implicita. Se l'ammissione locale
  viene rifiutata per pressione (profondit\`a o attesa stimata) e l'offload
  reattivo \`e abilitato, la richiesta viene deviata anzich\'e respinta.

  \item Il chiamante attende il futuro di completamento, con un proprio budget
  (\code{X-Timeout-Ms} o \code{timeoutMs} della funzione). Allo scadere risponde
  \code{408} soltanto \emph{quel} chiamante: l'esecuzione condivisa prosegue
  (\cref{cap:06-nucleo}).

  \item L'esito viene tradotto in risposta HTTP. Se la funzione ha dichiarato un
  proprio codice di stato, questo diventa il codice della risposta e viene
  segnalato dall'header \code{X-NanoFaaS-Function-Status}.
\end{enumerate}

L'ordine del punto dell'ammissione \`e la parte non ovvia. La sequenza \emph{coda sincrona
$\to$ coda asincrona $\to$ dispatch diretto} significa che il percorso sincrono
degrada in modo continuo al ridursi dei moduli presenti, senza che il codice
contenga un test sul nome dei moduli: ciascun ripiego \`e il default no-op del
punto di estensione precedente. Con un modulo di coda la scelta fra le due code
e l'ordine di servizio sono quelle del motore di scheduling composto
(\cref{cap:09-async-queue,cap:10-sync-queue}).

\subsection{Percorso asincrono}

Il percorso \code{:enqueue} \`e pi\`u corto e ha una differenza qualitativa: non
degrada. Se nessun modulo di coda fornisce l'accodamento asincrono, l'accodatore dichiara
di non supportarlo e l'API risponde \code{501 Not Implemented}. La scelta
\`e deliberata: un'invocazione asincrona eseguita in linea non sarebbe
un'invocazione asincrona degradata, sarebbe una risposta scorretta al contratto,
perch\'e il chiamante ha ricevuto un identificatore assumendo di poter
disconnettersi.

\section{Superficie dell'API}

L'API \`e piccola, e le sue tre famiglie di endpoint corrispondono alle tre
responsabilit\`a della piattaforma.

\begin{center}
\small
\begin{adjustbox}{max width=\textwidth}
\begin{tabular}{@{}llp{6.6cm}@{}}
\toprule
\textbf{Metodo} & \textbf{Percorso} & \textbf{Ruolo} \\
\midrule
\code{POST}   & \code{/v1/functions} & Registra una funzione. \\
\code{GET}    & \code{/v1/functions} & Elenca le funzioni. \\
\code{GET}    & \code{/v1/functions/\{name\}} & Ispeziona una funzione. \\
\code{PATCH}  & \code{/v1/functions/\{name\}} & Aggiorna parzialmente la specifica. \\
\code{DELETE} & \code{/v1/functions/\{name\}} & Rimuove funzione e risorse gestite. \\
\code{PUT/GET} & \code{/v1/functions/\{name\}/replicas} & Imposta / legge le repliche (\code{DEPLOYMENT}). \\
\midrule
\code{POST} & \code{/v1/functions/\{name\}:invoke}  & Invocazione sincrona. \\
\code{POST} & \code{/v1/functions/\{name\}:enqueue} & Invocazione asincrona. \\
\code{GET}  & \code{/v1/executions/\{id\}}          & Stato di un'esecuzione. \\
\midrule
\code{POST} & \code{/v1/internal/executions/\{id\}:complete} &
Callback di completamento dal runtime. \\
\midrule
\code{GET} & \code{/actuator/health} \emph{(:8081)}     & Sonde di liveness/readiness. \\
\code{GET} & \code{/actuator/prometheus} \emph{(:8081)} & Metriche. \\
\midrule
\code{GET/PATCH} & \code{/v1/admin/runtime-config} & Configurazione a caldo (\cref{cap:14-runtime-config}). \\
\bottomrule
\end{tabular}
\end{adjustbox}
\end{center}

La convenzione \code{:verbo} sul nome della risorsa (\code{:invoke},
\code{:enqueue}) segue lo stile delle API di Google Cloud per le azioni che non
sono operazioni CRUD sulla risorsa stessa. \`E preferibile a un sottopercorso
(\code{/functions/\{name\}/invoke}) perch\'e non suggerisce l'esistenza di una
sotto-risorsa che si possa elencare o cancellare.

\subsection{Trattamento degli header}

Il controller applica una regola che vale la pena isolare perch\'e riguarda la
sicurezza del contratto pi\`u che l'ergonomia:

\begin{lstlisting}[language=Java]
/**
 * Rebuild the request with the caller's real headers, lower-cased. The body's own
 * {@code headers} field is overwritten, never merged: a caller must not be able to
 * forge a header it did not send.
 */
\end{lstlisting}

Il corpo della richiesta contiene un campo \code{headers} che l'handler della
funzione vedr\`a. Se il control plane fondesse quel campo con gli header HTTP
reali, un chiamante potrebbe iniettare nel campo un header che non ha inviato,
e una funzione che decide in base agli header sarebbe ingannabile. Il campo
viene quindi \emph{sovrascritto}, mai unito. Dagli header propagati sono esclusi
quelli hop-by-hop e quelli che il control plane lega gi\`a a un parametro
dedicato.

\subsection{Mappatura degli errori}

I codici di stato che l'API produce sono pochi e ciascuno corrisponde a una
condizione precisa; il \cref{cap:A-api} li elenca per endpoint. Vale la pena
anticipare i meno ovvi. Ogni corpo di errore generato dal gestore globale ha i
campi in ordine fisso (\code{error}, \code{message}, poi \code{details} se
presente): una mappa non ordinata cambia l'ordine di iterazione da una JVM
all'altra, e un chiamante che confronta il corpo come testo lo vedrebbe cambiare.

\code{429 Too Many Requests} ha quattro origini distinguibili. Il limitatore di
frequenza risponde senza corpo, con \code{Retry-After: 1}. Il rifiuto della coda
sincrona porta \code{Retry-After} e \code{X-Queue-Reject-Reason} con valore
\code{depth}, \code{est\_wait} o \code{timeout}. Il superamento di una quota di
ammissione porta \code{Retry-After: 1} e un corpo
\code{\{"error":"invocation\_quota\_exceeded","resource":...\}}. La coda per funzione
piena, o l'assenza di capacit\`a nel dispatch diretto, risponde \code{429} senza
altro. Rendere distinguibile la causa \`e necessario perch\'e le condizioni
richiedono reazioni diverse dal chiamante: la piattaforma \`e globalmente satura,
lo \`e la singola funzione, o \`e una risorsa di memoria a esserlo.

\code{410 Gone} segnala la replica di una chiave il cui esito \`e stato espulso
per capacit\`a, con \code{X-Execution-Id}; \code{413} un corpo oltre il confine di
ingresso; \code{409} con codice \code{FUNCTION\_REMOVAL\_PENDING} una funzione la cui
rimozione ha lasciato risorse (\cref{cap:06-nucleo}); \code{408} lo scadere del
budget di un singolo chiamante sincrono.

\code{502 Bad Gateway} e \code{504 Gateway Timeout} sono emessi solo dal
percorso di offload, e segnalano che la richiesta \`e stata deviata verso
un'istanza remota che ha fallito o non ha risposto in tempo. Non esiste
fallback locale, e il \cref{cap:13-offload} argomenta perch\'e.

\section{Modello a thread}

Riassumendo i thread che il processo utilizza a regime. La colonna
\emph{Presenza} distingue ci\`o che esiste sempre da ci\`o che dipende dai moduli o
dalla configurazione.

\begin{center}
\small
\begin{adjustbox}{max width=\textwidth}
\begin{tabular}{@{}lll@{}}
\toprule
\textbf{Gruppo} & \textbf{Presenza} & \textbf{Ruolo} \\
\midrule
Event loop Netty & sempre; pari ai core & I/O HTTP in ingresso e in uscita \\
\code{boundedElastic} & sempre; elastico, limitato & Percorso \code{:enqueue}; invocazioni con chiave \\
\code{execution-expiry} & sempre; 1 (daemon) & Scadenza amministrativa dei record vivi \\
\code{nanofaas-core-retry-timer} & sempre; 1 (daemon) & Riprove ritardate del profilo diretto \\
Timer del cancello di prontezza & sempre; 1 (daemon) & Timeout dell'attesa dei deployment \\
\code{nanofaas-scheduler-engine} & con un modulo di coda; 1 & Motore di scheduling composto (\cref{cap:09-async-queue}) \\
Manutenzione dello stimatore & con \code{sync-queue}; 1 & Stima dell'attesa (\cref{cap:10-sync-queue}) \\
\code{nanofaas-internal-scaler} & con \code{autoscaler}; 1 & Campionamento e scalatura \\
\code{nanofaas-concurrency-governor} & con \code{concurrency-control}; 1 & Controllo della concorrenza \\
\code{nanofaas-admin-config-} & con \code{runtime-config}; 1 & Cambi di configurazione \\
\bottomrule
\end{tabular}
\end{adjustbox}
\end{center}

Il janitor periodico dello store, che la prima versione elencava, non c'\`e pi\`u:
la scadenza \`e affidata alle cache Caffeine (\cref{cap:06-nucleo}) e a un unico
timer di sola pianificazione.
Nella configurazione priva di moduli restano l'event loop, lo scheduler elastico
e i tre timer daemon: \`e ancora un numero ridotto di gruppi per un'intera
piattaforma FaaS, e la misura pi\`u compatta di che cosa significhi il vincolo di
minimalit\`a enunciato nel \cref{cap:03-requisiti-principi}. Va detto per\`o che
non tutti i gruppi sono stati verificati sul processo in esecuzione: l'elenco
proviene dal codice dei componenti, non da un dump dei thread.
